\scanpage{111}
% archive foliation: 110

Faisceaux $l$-adiques admissibles et motifs virtuels.

\deleted{1. Soit $S$ un schéma; $F$ un $\mathbb Q_l$}

1. On se fixe un nb premier $l$. Les schémas envisagés sont premiers à $l$. (\uline{le cas le plus intéressant est celui des schémas loc. de type fini sur $\mathbb Z$}.) Les faisceaux envisagés sont les $\mathbb Q_l$-faisceaux \inserted{constructibles}. On note $\mathbb Q_l(1)$ le faisceau de Tate $T_l(G_m)\otimes_{\mathbb Z_l}\mathbb Q_l$, et le twist par sa puissance tensorielle $n$-ième est noté $F\longmapsto F(n)$.

\margin{se borner pour commencer à ce cas!}

\margin{On travaille ds la catégorie des faisceaux $l$-adiques \unsure{constants tordus} \illegible{1}.}

2.1 Un faisceau $F$ sur $S$ est dit \uline{spécial} \inserted{\uline{de poids $i$}} (s'il est isomorphe à un \deleted{facteur direct} \inserted{\unsure{quotient constant tordu}} d'un \inserted{sous} faisceau \inserted{constant tordu d'un faisceau} de la forme $R^i f_*(\mathbb Q_{lX})$, où $f:X\longrightarrow S$ est projectif et lisse \inserted{de dim relative $i$ \unsure{ptt}}). On note que si $i<0$, cela implique $F=0$; si $i=0$, alors \deleted{\unsure{ce qui}} \deleted{\unsure{implique}} \inserted{équivaut à:} $F$ localement constant.$^{(*)}$

\margin{Définition \unsure{canulée|annulée}: dire qu'on se ramène à cette situation de \unsure{t.f.} sur $\unsure{\operatorname{Sp}|\operatorname{Spec}}\mathbb Z$.}

\margin{$(*)$ \deleted{\unsure{effectif}, \illegible{3} \unsure{étale} \illegible{2} poids $i$, \illegible{2} quotient \illegible{3} \unsure{non} \illegible{3} $\mathbb Q_l$ \illegible{2} \uline{\unsure{admissible}}!}}

\uline{Lemme} 2.\unsure{2|1}.

a) Stabilité de ``spécial de poids $i$'' \inserted{\unsure{sous-objets} et quotients cst tordus, et} par sommes directes.

b) Si $F,\unsure{F'}$ sont spéciaux de poids respectifs $i,j$, alors $F\otimes \unsure{F'}$ est spécial de $i+j$.

c) Stabilité par image inverse.

d) Si $F$ spécial de poids $i$, $\check F(-i)$ l'est aussi.

e) $\mathbb Q_l(-i)$ est spécial de poids $2i$ pour $i\geq 0$,

a) Se prouve en introduisant un $\coprod_{1\leq i\leq n}X_i$.

b) Se prouve en considérant un $X_1\times_S X_2$, et utilisant Künneth.

c) Par th. de changement de base.

d) \uline{Par th. de Lefschetz} en cohomologie $l$-adique.

\margin{d) on admet Lefschetz}

e) On \deleted{\unsure{prend}} a $\mathbb Q_l(-i)=R^{2i}f_*(\mathbb Q_{l\mathbb P^i_S})$.

\margin{\uline{Cor} 2.\unsure{3|2} Si $F$ spécial de poids $i$, et $n\geq i$, alors $\check F(-n)$ est spécial de poids $2n-i$}

3.1 Un faisceau $F$ sur $S$ est dit \uline{effectivement admissible} si \inserted{pour tout ouvert affine $U$ de $S$} l'on peut en trouver une \deleted{filtration \unsure{dont les quotients}} \inserted{filtration de $F|U$ dont les quotients} successifs sont de la forme $i_!(G)$, où $i:Y\longrightarrow \deleted{X}\inserted{U}$ est une immers. de présentation finie et $G$ sur \deleted{$\unsure{S}$} $Y$ est \deleted{\unsure{sont de la forme} \unsure{image directe} \unsure{des} \unsure{ouverts de} $U$ \unsure{des parties}} \deleted{\unsure{des faisceaux} \unsure{admissibles} \unsure{les} $F|S_{\unsure{\alpha}}$ \unsure{sont}} -admissible.

\margin{Il faut se permettre une extension \unsure{finie surj. loc.} \unsure{radicielle} de $Y$, pour rendre $G$ spécial}

\deleted{On dit} Si $\Phi$ est une partie de $\mathbb Z$, on dit que $F$ est à poids $\subset\Phi$ si \deleted{les $F|S_{\unsure{\alpha}}$ sont} \inserted{on peut prendre les $G$} admissibles de poids $\in\Phi$. (\uline{De eff admissible} alors $F$ est de poids $\subset\Phi'$, pour tt $\Phi'\supset\Phi$). Si \deleted{$\unsure{\Phi\subset E}$ est} \deleted{$\Phi$, cela implique} $\Phi$ ne contient pas d'entiers $\geq0$, cela implique $F=0$.

\margin{Attention, il n'est pas tautologique que si $F$ est eff! adm. à poids $\in\Phi$ et à poids $\subset\Phi'$, \uline{il soit} à \deleted{poids} \inserted{poids} $\subset\Phi\cap\Phi'$, \unsure{même} $\Phi\cap\Phi'=\varnothing$.}
